% ===========================================================================
% Section 6: Mode-Specific Patterns via Counterfactuals
%
% Robot-mode (Shapiro Step 3): state facts, do not argue; define concepts
% before use; no forward references to undefined constructs.
%
% Anchors: Table 13 (tab:counterfactual), Figure C13 (fig:ma_counterfactual).
%
% PROVISIONAL FRAMING (decisions log F1): "both modes contribute, the rail
% component is better identified." The alternative ("rail dominates") is a
% one-subsection rewrite; see Plan/provisional_decisions_log.md.
%
% Caveat language per task A3: "suggestive evidence," "characterize
% relative importance," never "the effect of rail closures is X."
%
% Numbers inline for now, from results/tables/table_13_counterfactual.csv;
% AutoFill macros in the paper-wide substitution pass.
%
% CITATION KEYS: none new beyond Sections 1-5.
% ===========================================================================

\section{Mode-Specific Patterns via Counterfactuals}
\label{sec:counterfactuals}

The estimates in Section~\ref{sec:results} measure the effect of the
combined change in market access produced by the simultaneous
contraction of the rail network and expansion of the road network.
This section decomposes that combined change into its two modal
components. The decomposition holds one network fixed at its
pre-reform configuration while the other changes as observed, and
recomputes market access under each hypothetical configuration.

\subsection{Counterfactual Market Access Measures}
\label{sec:cf_construction}

We construct two counterfactual market access measures using the same
cost-surface pipeline as the actual measures described in
Section~\ref{sec:data}. The \emph{rail-only} measure,
$\Delta \ln \mathrm{MA}^{\mathrm{rail}}_i$, combines the 1986 rail
network with the road network frozen at its 1954 configuration; it
isolates the market-access consequences of the rail contraction. The
\emph{road-only} measure, $\Delta \ln \mathrm{MA}^{\mathrm{road}}_i$,
combines the road network of 1986 with the rail network frozen at its
1960 configuration; it isolates the consequences of the road
expansion. Both measures use the same trade elasticity
($\theta = 4.55$), the same destination populations, and the same
navigation layer as the full measure, so differences across the three
measures reflect only the network configuration.

Each counterfactual regressor is instrumented with the instrument that
targets its source of variation. The rail-only measure is instrumented
with the Larkin Plan instrument alone: the Plan's study designations
shift rail-driven market access and have no mechanical connection to
the frozen road layer. The road-only measure is instrumented with the
hypothetical-road instrument alone. The full measure, reproduced from
Table~\ref{tab:population_iv} for comparison, uses both instruments.
The three regressors are never entered jointly in one regression: the
full change is by construction close to the sum of its components, and
joint estimation is not informative in a sample of 309 districts.

Figure~\ref{fig:ma_counterfactual} maps the three measures on a common
scale. The rail-only change is negative in every district: the median
district loses 3~log points of market access from the rail contraction
alone (interquartile range $-0.09$ to $-0.02$), and a small number of
remote districts lose substantially more. The road-only change is
positive for nearly all districts and reproduces the broad spatial
pattern of the full measure, consistent with the road expansion being
the dominant source of the aggregate market-access gains documented in
Section~\ref{sec:data}.

\subsection{Results}
\label{sec:cf_results}

Table~\ref{tab:counterfactual} reports estimates of
equation~\eqref{eq:main} with each of the three regressors, for the
four population outcomes of Table~\ref{tab:population_iv}. Panel~A
reproduces the full-measure estimates. Panels~B and~C report the
rail-only and road-only estimates.

For total population, the IV estimate on the rail-only measure is
0.032 (0.034) with a first-stage $F$ of 105.0. The IV estimate on the
road-only measure is 0.039 (0.051) with a first-stage $F$ of 13.2.
The full-measure estimate is 0.046 (0.033) with an $F$ of 13.6. The
three point estimates are of similar magnitude; none is statistically
distinguishable from zero at conventional levels, and none is
distinguishable from the others. What differs sharply across panels is
identification strength: the Larkin Plan instrument predicts
rail-driven market access far more strongly ($F$ between 69.6 and
117.8 across outcomes) than either instrument predicts its respective
measure in the other panels.

The urban and rural decomposition shows the same pattern. Rail-only
estimates are 0.045 (0.035) for urban and 0.078 (0.084) for rural
population; road-only estimates are 0.026 (0.067) and 0.098 (0.133).
The urban-share estimates are near zero in all three panels, as in
Table~\ref{tab:population_iv}. OLS estimates differ across panels: the
rail-only OLS coefficient is $-0.004$ (0.019), while IV moves it to
0.032, a pattern consistent with rail lines having been closed in
districts that were already declining.

\subsection{Interpretation and Caveats}
\label{sec:cf_caveats}

Three features of this exercise limit what the estimates can support,
and we state them before drawing any conclusion from the comparison.

First, the counterfactual networks are accounting constructs, not
equilibrium predictions. Freezing the road network at 1954 does not
model how firms, households, or the government would have behaved had
the road program not occurred. The estimates therefore characterize
the relative importance of the two modal shocks for the market-access
channel; they are not estimates of the effect of a counterfactual
policy.

Second, the exclusion restrictions are inherited from
Section~\ref{sec:identification} and are not separately testable here. The
Larkin Plan instrument is excludable for the rail-only measure under
the same conditions as for the full measure. For the road-only
measure, the hypothetical-network instrument's first-stage $F$ of 13.2
for total population falls to 9.1 for urban population, so the
road-only panel is less precisely identified throughout.

Third, differences in precision across panels reflect the instruments,
not the economics. The rail-only panel is the best-identified object
in this paper ($F = 105.0$); this makes its estimate the most
credible, not the largest or the most important.

With these limits stated, the decomposition supports one suggestive
conclusion: both modal shocks contribute to the population response,
with similar point estimates, and the rail component is identified
with far greater precision. The point estimates give no support to the
view that one mode's market-access channel dominates the other; the
data are instead consistent with population responding to market
access itself, whichever mode produces it. Section~\ref{sec:mechanisms}
examines how much of this response operates through infrastructure
located in the district itself.
